Q01 & Adhérence du SI aux missions & \textbf{3} & P01, P03, P04 & MEDICA-SERVICES vend l'hébergement et l'exploitation d'applications pour 12 établissements de soins (P01). Sans SI il n'existe plus de prestation et les clients reviennent au papier (P04). Le SI est vital aux missions. & Confirmé \\
Q02 & Impacts internes & \textbf{2} & P01, P04 & La DG prévoit de fortes pertes contractuelles après une journée d'arrêt et le départ de certains clients si l'arrêt se prolonge (P04). Elle estime pouvoir poursuivre des prestations manuelles, ce qui écarte une conséquence fatale démontrée. & Probable \\
Q03 & Impacts externes & \textbf{3} & P04, P05 & Les établissements clients basculent sur papier (P04) et une altération non détectée peut provoquer des erreurs d'orientation et de prise en charge de patients (P05). La sécurité de personnes chez des tiers est en jeu. Note fondée sur des déclarations, sans étude d'impact chiffrée. & Probable \\
Q04 & Disponibilité & \textbf{2} & P01, P04, P05 & Une journée d'arrêt du portail est grave (pertes contractuelles, P04) et un arrêt prolongé affecte fortement la continuité des soins (P05). La DG décrit une activité dégradée et non un arrêt définitif, d'où 2 et non 3. Reprise contractuelle sous 8 h (P01). & Probable \\
Q05 & Intégrité & \textbf{2} & P05 & Les environnements contiennent des identifiants patients et des informations de santé. Une altération non détectée peut provoquer des erreurs d'orientation et de prise en charge (P05). L'impact est grave pour l'activité. Aucune pièce ne le décrit comme fatal pour l'entreprise, d'où 2 par cohérence avec Q02 et Q04. & Probable \\
Q06 & Confidentialité & \textbf{3} & P05, P07 & Identifiants patients et données de santé (P05). La DSI qualifie la divulgation de risque critique et les attaquants visés cherchent la revente de données et l'extorsion (P07). & Probable \\
Q07 & Fréquence des sinistres & \textbf{2} & P06 & Dix événements SSI enregistrés en 12 mois, soit environ 2,5 par trimestre (P06). Le registre n'étant ni exhaustif ni revu, la fréquence réelle peut être supérieure. & Probable \\
Q08 & Motivation des attaquants & \textbf{2} & P07 & Menaces ciblées jugées plausibles et potentiellement fortement motivées par la revente de données et l'extorsion (P07). Aucune pièce ne justifie le niveau très important. & Probable \\
Q09 & Moyens des attaquants & \textbf{2} & P07 & P07 écarte les moyens illimités. Une cybercriminalité d'extorsion dispose plausiblement de moyens importants. Hypothèse d'auditeur sans effet sur le maximum du groupe, déjà fixé à 2 par Q07 et Q08. & Incertain \\
Q10 & Hétérogénéité du SI & \textbf{2} & P08, T01 & Windows, Linux, hyperviseurs, réseau multi-constructeurs, SaaS et application propriétaire héritée. P08 qualifie l'architecture de fortement hétérogène mais non extrêmement hétérogène. & Confirmé \\
Q11 & Ouverture du SI & \textbf{2} & P03, P07, T01, T02 & Interconnexions clients encadrées par contrat (P07). Le portail est ouvert à des internautes non maîtrisés, mais par une publication unique filtrée (FW-01, FW-08) et un reverse proxy en DMZ (P03, T01, T02). Ouverture externe jugée sous contrôle. La note 3 serait défendable si l'on considère que l'accès public au portail échappe au contrôle de l'organisme. La somme passerait alors à 11 et la cible à 5. & Probable \\
Q12 & Variabilité du SI & \textbf{2} & P06, P08 & Changements applicatifs mensuels et d'infrastructure trimestriels (P08), deux indisponibilités dues à des changements (P06). Le SI change relativement souvent. & Confirmé \\
